\documentclass{article}

\usepackage[utf8]{inputenc}     % pdfLaTeX
\usepackage[T1]{fontenc}
\usepackage[magyar]{babel}
\usepackage{amsmath}

% Set page size and margins
% Replace `letterpaper' with `a4paper' for UK/EU standard size
\usepackage[a4paper,top=2cm,bottom=2cm,left=3cm,right=3cm,marginparwidth=1.75cm]{geometry}
%for arrows
\usepackage{textcomp}

\title{Névelők}
\author{Lengyel Zoltán Péter\thanks{Csukás Ibolya tanítása alapján}}
\date{Legutóbb frissítve: 2026. Március 20.}

\begin{document}

\maketitle
\tableofcontents
\newpage

\begin{center}
    \section{Névelők}
\end{center}

\begin{itemize}
    \item Főnevek vagy főnévi szerkezetek előtt állnak 
    \item Utal a főnév határozott vagy határozatlan voltára
    \item Határozott névelők: a, az
    \item Határozatlan névelő: egy\footnote{Az 'egy' számnévből ered, hangsúlytalanná vált}
\end{itemize}

\end{document}